%% ============================================================
%%  SECCIÓN 8 — Análisis Experimental
%% ============================================================

\section{Análisis Experimental}

En esta sección se examina el comportamiento cualitativo de los materiales
a partir de los datos experimentales presentados en la
Sección~\ref{sec:resultados}. El análisis se limita a la identificación de
tendencias y patrones observables, sin considerar aún el cálculo de
magnitudes derivadas.

%% ------------------------------------------------------------------
\subsection{Comportamiento del Resorte Metálico}
%% ------------------------------------------------------------------

De acuerdo con la Tabla~\ref{tab:datos_resorte}, la longitud $L$ del
resorte presenta un incremento monótono al aumentar la masa aplicada,
variando desde \qty{29,8}{cm} hasta \qty{44,6}{cm}. Los incrementos
sucesivos de longitud muestran una tendencia aproximadamente uniforme,
lo que sugiere una relación proporcional entre la carga y la deformación
longitudinal.

En relación con el diámetro $D$, se observa una disminución progresiva
conforme aumenta la carga, pasando de \qty{14,22}{mm} a
\qty{13,90}{mm}. Esta variación es consistente con la contracción
transversal asociada a la elongación del material.

Durante la fase de descarga, las longitudes registradas para cada
configuración de masa coincidieron con los valores obtenidos en la fase
de carga. Al retirar todas las masas, el resorte recuperó íntegramente
su longitud natural, sin evidencia de deformación remanente. Este
comportamiento confirma que el resorte permaneció dentro de su región
elástica durante todo el experimento, por lo que no se presenta una
tabla de descarga separada.

%% ------------------------------------------------------------------
\subsection{Comportamiento de la Liga Elástica}
%% ------------------------------------------------------------------

\subsubsection{Fase de carga}

Durante la fase de carga (Tabla~\ref{tab:datos_liga}), la longitud $L$
de la liga aumenta con la masa aplicada, desde \qty{36,7}{cm} hasta
\qty{48,3}{cm}. A diferencia del resorte, los incrementos de longitud
no presentan una tendencia uniforme, evidenciando una relación no lineal
entre la carga y la deformación.

Las dimensiones de la sección transversal $a$ y $b$ disminuyen conforme
se incrementa la carga. No obstante, en ciertas mediciones consecutivas
se observa una variación mínima o nula, lo que podría estar asociado a
limitaciones experimentales o a cambios en la respuesta mecánica del
material en ese intervalo.

\subsubsection{Fase de descarga}

En la fase de descarga, al comparar valores correspondientes a una misma
masa, se observa que la longitud $L$ es sistemáticamente mayor que en la
fase de carga. Por ejemplo, para \qty{1000}{g} se registra una longitud
de \qty{41,5}{cm} en carga y \qty{45,5}{cm} en descarga, una diferencia
de \qty{4,0}{cm}.

Esta diferencia es considerable y puede atribuirse a tres mecanismos no
excluyentes entre sí:

\begin{enumerate}
    \item \textbf{Histéresis viscoelástica:} En elastómeros, las cadenas
    poliméricas no recuperan su configuración de equilibrio de forma
    inmediata al disminuir la carga, produciendo una recuperación retardada
    que genera longitudes de descarga mayores que las de carga para igual
    masa.

    \item \textbf{Deformación plástica parcial:} Si la liga superó
    localmente su límite elástico durante el ensayo, parte de la deformación
    se vuelve permanente e irreversible, contribuyendo a la diferencia
    observada.

    \item \textbf{Errores sistemáticos de medición:} La resolución limitada
    de la regla métrica (\qty{\pm 0,25}{cm} para la liga) y el posible
    error de paralaje pudieron introducir desviaciones sistemáticas,
    especialmente en las mediciones de descarga donde el material aún no
    había alcanzado el equilibrio estático.
\end{enumerate}

De manera similar, las dimensiones transversales $a$ y $b$ no coinciden
con los valores de carga para iguales condiciones de masa, con valores
de descarga sistemáticamente menores que los de carga en la región de
cargas altas. Este comportamiento en conjunto es característico de
materiales \textbf{viscoelásticos} con posible componente plástica,
cuya respuesta mecánica depende de la historia de carga y no es
completamente reversible.